\documentclass[10pt]{article}
\usepackage{cmap}
\usepackage[T1]{fontenc}
\usepackage{lmodern}
\usepackage{multicol}
\usepackage{fancyhdr}
\usepackage[bookmarks=true]{hyperref}
\hypersetup{pdftitle={Sample Book}, pdfauthor={Ada Example}}
\setcounter{secnumdepth}{0}
\setlength{\columnsep}{30pt}
\pagestyle{fancy}
\fancyhf{}
\fancyhead[C]{SAMPLE BOOK}
\fancyfoot[C]{\thepage}
\renewcommand{\headrulewidth}{0pt}

\begin{document}
\pagenumbering{roman}
\thispagestyle{empty}
\vspace*{6\baselineskip}
\begin{center}
{\huge\bfseries Sample Book\par}
\vspace{2\baselineskip}
{\large\bfseries Ada Example\par}
\end{center}
\vspace{4\baselineskip}

This small book is a test fixture. It has a title page numbered in roman
numerals, four body pages numbered from one, a running header, and a page
where the text is set in two columns. Nothing in it is meant to be read for
its own sake, but every sentence is complete so that paragraphs end where a
reader expects them to end.

\clearpage
\pagenumbering{arabic}

\section{Entropy}
Entropy measures how much surprise a source of messages carries. A fair coin
has one bit of entropy, because each toss is equally likely to land either
way, while a coin that always lands heads has none, because nothing it says
can surprise anyone.

Every message that a source can send carries some amount of infor-\linebreak
mation, and the content of one message is the logarithm of the inverse of its
probability. Rare messages carry more of it, and the average over all messages
is the entropy of the source. This is what makes entropy the natural measure
of how far a message can be compressed.

A claim that is self-evident carries no surprise to those who already hold
it, and so it adds no information for them. The same message can be news to
one reader and noise to another, which is why entropy is always stated
relative to what the receiver knows. In 1948 the idea was set out in full,
and the figure 1996 below is the number of a later survey, not a year.

\subsection{Joint entropy}
Joint entropy is the entropy of two variables taken together. If the two are
independent, it is the sum of their separate entropies. If one tells you
something about the other, the joint entropy is smaller than that sum, and
the difference is the mutual information between them.

Consider two coins tossed by the same hand. When the first lands heads the
second is a little more likely to do so as well, since the hand repeats its
own motion. The pair of results then carries less than two full bits, and a
code for the pair can use fewer than two bits on average.

\newpage

Conditional entropy is the entropy that remains in one variable once the
other is known. It never exceeds the entropy of the variable on its own, and
it equals that entropy only when the two variables are independent. A weather
report that always agrees with the weather leaves no conditional entropy at
all, while a report written at random leaves all of it.

The chain rule ties these quantities together. The joint entropy of two
variables is the entropy of the first plus the conditional entropy of the
second given the first. Applied again and again, the rule splits the entropy
of a long message into the cost of each symbol given the symbols before it.

Students are told early in the course that a thing which is self-\linebreak
evident to a theorist may still need a proof, and the proof of the chain rule
is short enough to give in full. It follows from the definition of conditional
probability and the additivity of the logarithm, and it needs nothing else.

Mutual information is the difference between an entropy and a conditional
entropy. It is symmetric, it is never negative, and it is zero exactly when
the two variables are independent. Every later result in this book about
what a channel can carry is a statement about this one quantity.

\newpage

\begin{multicols}{2}
Alpha column opens the page with the first half of the text. Take a source
that sends the letter A with probability one half and the letters B and C
with probability one quarter each. Its entropy is one and a half bits, which
is less than the logarithm of three that a source with three equally likely
letters would have. A code that spends one bit on A and two bits on each of B
and C reaches that figure exactly, and no code can do better on average.
\columnbreak

Omega column holds the second half of the text. Now let two such sources run
side by side without influencing one another. Their joint entropy is three
bits, the sum of the two. Link them, so that the second repeats the first one
time in four, and the joint entropy falls below three bits by exactly the
information that each now carries about the other. The saving is real, and a
joint code can collect it.
\end{multicols}

\newpage

\section{Channels}
A channel takes the symbols a sender puts in and returns the symbols a
receiver gets out. Noise makes the two differ. The capacity of a channel is
the highest rate, in bits per symbol, at which messages can be sent so that
the receiver recovers them with an error as small as one likes.

The noisy channel coding theorem says that every rate below the capacity can
be reached and that no rate above it can. The proof is not constructive, and
for decades it was not known how to build codes that come close to the limit.
Modern codes do come close, and the gap between theory and practice is now
small enough to measure in fractions of a decibel.

The theorem ends this short book. Entropy told us how much a source says,
mutual information told us how much a channel passes on, and the capacity
told us the most that any code could ever ask of it. Everything else is the
engineering of getting there.
\end{document}
